\section{Sharper clamped and material estimates}\label{inf:sec-sharp}

The spaces in this section remain $\Cc$ and $\Xc_p=\Gc_p\times\Ec_1$.
The centre and the fixed material field have the membership supplied
by Section~\ref{inf:sec-material}. The estimates below follow from
their exact equations.
For fixed wider parameter ranges use the extensions in
Lemmas~\ref{inf:split-extension} and \ref{inf:fixed-detuning-extension}.

\begin{lemma}[The harmonic gradient coefficient sum]\label{inf:pole-sum}
For endpoint-normalized solid harmonics $S_i,S_j$,
\begin{equation}\label{inf:pole-sum-identity}
 \nabla S_i\cdot\nabla S_j
 =\sum_{k<i+j}2h(h+2k+p-1)c_{ij}^kS_kt^{h-1},\quad h=i+j-k,
 \qquad
 \sum_{k<i+j}2h(h+2k+p-1)c_{ij}^k=4ij.
\end{equation}
Consequently, for $H_u=\HP h_u$,
\begin{equation}\label{inf:sharp-gradient-multiplier}
 \nrm{M_{\nabla H_1\cdot\nabla H_2}}_{\Cc\to\Cc}
 \le4\nrm{h_1}_{\Ec_1}\nrm{h_2}_{\Ec_1}.
\end{equation}
\end{lemma}
\begin{proof}
The product rule for two harmonic polynomials gives
$2\nabla S_i\cdot\nabla S_j=\Delta(S_iS_j)$.
Apply \eqref{inf:profile-laplacian} to each term of the positive solid
product to obtain the first identity. At a unit point of the $x=1$
block axis, $S_i=S_j=1$ and their tangential gradients vanish.
Their full gradients are $2iy$ and $2jy$. Evaluation therefore gives
the second identity; the coefficients are all nonnegative.

Every multiplier $S_kt^{h-1}$ has total degree $2i+2j-2$ and norm
at most $2^k\wt(2i+2j-2)$ by Proposition~\ref{inf:module}.
This is at most $2^{i+j}\wt(2i)\wt(2j)$.
The positive coefficient sum is $4ij\le4(1+i)(1+j)$.
Sum the absolute shape coefficients with their $\Ec_1$ weights
to get \eqref{inf:sharp-gradient-multiplier}. A constant harmonic
has zero gradient, and completion gives the assertion for all shapes.
\end{proof}

\begin{proposition}[Sharp centre and material bounds]\label{inf:sharp-material}
Under the root and centre hypotheses of
Lemma~\ref{inf:detuning-dependencies}, for each fixed centre order
$N\ge4$, uniformly in
$0<\tau\le T\le\log p$ and on the working region
\eqref{inf:nonlinear-working-region},
\begin{equation}\label{inf:sharp-three-bounds}
 \nrm{g_0}_\Cc\le p^3\mathcal B_N(T),\qquad
 \nrm{\Qmat^{-1}}\le p^2\mathcal B_N(T),\qquad
 \sup\nrm{D^2\Fnl}\le p^3\mathcal B_N(T).
\end{equation}
Here $\mathcal B_N(T)=C_N(1+T)^{d_N^{\mathrm{tame}}}e^{C_N\sqrt T}$
is the envelope of Lemma~\ref{inf:detuning-dependencies}, enlarged
if necessary; $C_N,d_N^{\mathrm{tame}}$ depend on the fixed order
and compact split interval. For fixed $T$ it is a constant.
The inverse identities are those of
Theorem~\ref{inf:material-isomorphism}.
\end{proposition}
\begin{proof}
In the proof of Proposition~\ref{inf:nonlinear-map}, replace
\eqref{inf:gradient-multiplier} by
\eqref{inf:sharp-gradient-multiplier}. Every other grouped estimate
there has a constant independent of $p$.
In particular \eqref{inf:second-shape-clamped} and
\eqref{inf:mixed-shape-clamped} are applied to clamped sources, as
their statements require. The inverse multiplier series and their
first two derivatives are bounded on the working ball. Thus
\begin{equation}\label{inf:sharp-nonlinear}
 \nrm{D_f\Fnl(g,f)}\le C\nrm g_\Cc,
 \qquad \nrm{D^2\Fnl(g,f)}\le C(1+\nrm g_\Cc).
\end{equation}
The frequency term is affine in $g$ and independent of the shape.

Put $\mathcal R_{\mathrm{sc}}=(P_{\Phi_N}\Delta-\Delta)\KD$.
Equation~\eqref{inf:relative-scalar} gives
$\nrm{\mathcal R_{\mathrm{sc}}}\le C\tau p^{-2}<1/2$ at the fixed
centre. Since $U_0=1+Kg_0$, its residual equation is exactly
\[
 E_0=(I+\mathcal R_{\mathrm{sc}})g_0+R^2U_0.
\]
Use $\nrm{g_0}\le2\nrm{(I+\mathcal R_{\mathrm{sc}})g_0}$ and
$\nrm{U_0}\le C_Np$ from \eqref{inf:centre-norm-bounds}.
This yields $\nrm{g_0}\le C_Np^3$ and hence the last bound in
\eqref{inf:sharp-three-bounds} by \eqref{inf:sharp-nonlinear}.
This is norm absorption on the already defined clamped source.

The fixed material field is
$\mathfrak f_0=d_0^{-1}\Euler U_0=d_0^{-1}\Euler Kg_0$.
The all-source estimate for $\Euler\KD$ in
\eqref{inf:poisson-operator-bounds} gives
$\nrm{\mathfrak f_0}\le C_Np^2$. Its zero Dirichlet trace and
$\Delta\mathfrak f_0\in\Cc$, already proved in
\eqref{inf:fixed-material-field}, give
$\mathfrak f_0=\KD\Delta\mathfrak f_0$.
Apply the material identity to the constant shape variation one:
\[
 B_{\Phi_N}\mathfrak f_0
 =d_0^{-1}\Euler E_0-D_f\Fnl(z_0)[1].
\]
The residual bound and \eqref{inf:sharp-nonlinear} bound this by
$C_Np^3$. On the other hand it equals
$(I+\mathcal R_{\mathrm{sc}})\Delta\mathfrak f_0+R^2\mathfrak f_0$.
The same norm absorption proves
\begin{equation}\label{inf:sharp-material-field}
 \nrm{\Delta\mathfrak f_0}\le C_Np^4.
\end{equation}
For an unknown shape $\eta$,
\[
 \Delta T_\eta=(\HP\eta)\Delta\mathfrak f_0
 +2\nabla(\HP\eta)\cdot\nabla\KD\Delta\mathfrak f_0.
\]
The module estimate and the all-source gradient bound
\eqref{inf:gradient-C-one} give
$\nrm{\Delta T_\eta}\le C_Np^4\nrm\eta_{\Ec_1}$.
The exact reconstruction \eqref{inf:material-reconstruction} has
$\nrm\eta_{\Ec_1}\le C_Np^{-2}\nrm\psi_\Cc$, by its normal trace
and the normal Hessian bound. It follows that
$\nrm{\delta g}+\nrm\eta\le C_Np^2\nrm\psi$.
Both inverse identities were proved before estimating this expression,
so this proves the middle bound in \eqref{inf:sharp-three-bounds}.
The dependence on $T$ in these same steps is given by row 12 and
the clamped-source and material calculations in the proof of
Lemma~\ref{inf:detuning-dependencies}, yielding the stated envelope.
\end{proof}

\begin{lemma}[Quadratic graph remainder at a polynomial chart]
\label{inf:quadratic-graph}
Let $H$ be a real invariant harmonic polynomial of fixed angular degree,
with size $\delta$ in a fixed multiplier norm including its first
four Euler derivatives. For $p\delta$ sufficiently small put
$q=1+H$, $d=q+\Euler q$, $m=q^{p-1/2}d^{1/2}$, and denote its
unitary Dirichlet operator by $\widehat H_H$. Its first variation is
\[
 V_1(H)=2H\Delta+2\nabla H\cdot\nabla\Euler
                         +2\nabla l_{\mathrm{lin}}\cdot\nabla,
 \qquad l_{\mathrm{lin}}=pH+\tfrac12\Euler H.
\]
In each of the mixed and Hilbert matrix norms at fixed exponential
weight,
\begin{equation}\label{inf:quadratic-graph-bound}
 \nrm{(\widehat H_H-H_0-V_1(H))K_0^{-1}}\le C\delta^2.
\end{equation}
The constant may depend on the fixed degree and weights. This is
an all-source estimate at a fixed polynomial chart.
\end{lemma}
\begin{proof}
Lemma~\ref{inf:pole-sum} also gives a $p$-independent multiplier
bound for fixed-degree gradient products in $\Mmix$: each positive
solid multiplier is bounded by Lemma~\ref{inf:mixed-multipliers}.
In $\Hc$ use its fixed bandwidth and the pointwise gradient bound
as in \eqref{inf:Hilbert-gradient-product}. These bounds apply to
$H$ and each of its fixed Euler derivatives.
With $l=\log m$, harmonicity of $q,d$ gives
\[
 \Delta l=-(p-\tfrac12)q^{-2}|\nabla H|^2
       -\tfrac12d^{-2}|\nabla(1+\Euler)H|^2.
\]
Thus $\Delta l=O(p\delta^2)$,
$|\nabla l|^2=O(p^2\delta^2)$, and
$\Euler l,\Euler^2l=O(p\delta)$ in the multiplier norms.
Subtracting $l_{\mathrm{lin}}$ in its gradient leaves harmonic
gradients times coefficients of size $O(p\delta)$; after
$\nabla\KD$ the result is $O(\delta^2)$.
The exact conjugations are
\begin{align*}
 m\Delta m^{-1}&=\Delta-2\nabla l\cdot\nabla
                            +|\nabla l|^2-\Delta l,\\
 m(\nabla H\cdot\nabla\Euler)m^{-1}
 &=\nabla H\cdot\nabla\Euler-(\nabla H\cdot\nabla l)\Euler
       -(\Euler l)\nabla H\cdot\nabla\\
 &\quad +(\Euler l)(\nabla H\cdot\nabla l)
                         -\nabla H\cdot\nabla\Euler l,\\
 m\Euler(\Euler+1)m^{-1}
 &=\Euler(\Euler+1)-2(\Euler l)\Euler
             +(\Euler l)^2-\Euler^2l-\Euler l,\\
 m\Euler m^{-1}&=\Euler-\Euler l.
\end{align*}
Insert these in the negative scalar pullback
\eqref{inf:scalar-pullback}. Its linear part is exactly $V_1(H)$;
there is no linear potential because $\Delta l_{\mathrm{lin}}=0$.
The following table records every remainder after composition with
$\KD$; all Poisson and gradient bounds are the all-layer bounds in
Lemma~\ref{inf:mixed-Poisson}.
\begin{center}\small\normalfont
\begin{tabular}{@{}>{\raggedright\arraybackslash}p{.60\textwidth}>{\raggedright\arraybackslash}p{.33\textwidth}@{}}
\toprule Remainder & Bound after $\KD$\\\midrule
Quadratic part of $q^{-2}\Delta$ & $C\delta^2$\\
Nonlinear density-gradient terms & $Cp\delta^2\cdot p^{-1}$\\
Density potential & $Cp^2\delta^2\cdot p^{-2}$\\
Nonlinear coefficient of $\nabla H\cdot\nabla\Euler$
 & $C\delta\cdot\delta$\\
$(\nabla H\cdot\nabla l)\Euler$ and
$(\Euler l)\nabla H\cdot\nabla$ & $Cp\delta^2\cdot p^{-1}$\\
Their zero-order terms & $C(p^2\delta^3+p\delta^2)p^{-2}$\\
$|\nabla H|^2\Euler(\Euler+1)$ & $C\delta^2$\\
Its density corrections & $C(\delta^3+\delta^4+\delta^3/p)$\\
The scalar $\Euler$ coefficient involving $C_q$ & $C\delta^2/p$\\
Its density potential & $Cp\delta^3\cdot p^{-2}$\\
\bottomrule
\end{tabular}
\end{center}
Formula~\eqref{inf:Cq-identity} makes $C_q$ quadratic in the fixed
harmonic gradients and Euler derivatives. The inverse multipliers
have convergent series in both matrix algebras. Every row is bounded
by $C\delta^2$, including the possible $p^2$ potentials. The identities
first hold on the polynomial source core and extend by the graph
bounds. Multiplying the $\KD$ version by the angular-diagonal
contraction $H_0K_0^{-1}$ proves \eqref{inf:quadratic-graph-bound}.
\end{proof}
